
\subsection{Qualitative Case Study: Resolution of Modality Conflict}
\label{subsec:case_study}
To demonstrate the framework's capability in handling discordant evidence, we analyze a query exhibiting a deliberate semantic conflict between visual and textual modalities.

\textbf{Query Input:}
\begin{itemize}
    \item \textbf{Text:} ``Hi I hit the side of my forehead... then the lump went down and the side of my forehead swelled up... now my eyelid is black and swollen...'' (Describes physical trauma/hematoma).
    \item \textbf{Image:} A dermoscopic image of a \textit{Skin Growth} (seborrheic keratosis).
\end{itemize}

Arguments from both modalities are clinically distinct: the text strongly indicates a trauma-induced specific condition (\textit{Swollen Eye}), while the visual evidence supports a chronic dermatological condition (\textit{Skin Growth}).

\textbf{System Response Analysis:}
Table~\ref{tab:case_study} details the counterfactual belief distributions. The framework's behavior reveals three critical safety mechanisms in action:

\begin{enumerate}
    \item \textbf{Modality Disentanglement:} The \textit{No Image} counterfactual correctly identifies \textit{Swollen Eye} ($p=0.85$) as the primary diagnosis based on narrative symptoms. Conversely, the \textit{No Text} counterfactual pivots entirely to \textit{Skin Growth} ($p=1.0$), confirming that the visual encoder successfully recognizes the pathology despite the distracting narrative.
    
    \item \textbf{Conflict Detection:} The \textit{Baseline} distribution reflects this ambiguity, splitting probability mass between \textit{Skin Growth} (0.64) and \textit{Swollen Eye} (0.26). Crucially, the \textbf{Modality Consistency} constraint score is low ($0.26$), explicitly signaling that the two information sources contradict each other.
    
    \item \textbf{Safe Failure:} Instead of forcing a high-confidence prediction, the system triggers a ``Low Robustness'' warning (Stability Divergence $= 0.39$) and flags the result as \textit{Near Decision Boundary}.
\end{enumerate}

This validates that CountRAG-Clinic effectively surfaces diagnostic ambiguity rather than suppressing it, preventing potential misdiagnosis in scenarios where patient-provided history contradicts visual findings.

\begin{table}[h]
\caption{Counterfactual Belief Shifts under Modality Conflict}
\label{tab:case_study}
\centering
\resizebox{\columnwidth}{!}{%
\begin{tabular}{lccc}
\toprule
\textbf{Diagnosis} & \textbf{Baseline} & \textbf{No Text (Image Only)} & \textbf{No Image (Text Only)} \\ 
\midrule
Skin Growth & \textbf{0.64} & \textbf{1.00} & - \\
Swollen Eye & 0.26 & - & \textbf{0.85} \\
Dry Scalp & 0.10 & - & - \\
Itchy Eyelid & - & - & 0.15 \\
\midrule
\textit{Dominant Signal} & \textit{Conflict} & \textit{Visual} & \textit{Textual} \\
\bottomrule
\end{tabular}%
}
\end{table}
